% ============================================================
% Basic LaTeX Article Template — paper-draft skill default
% ============================================================
% Check CLAUDE.md for a project-specific template before using this one.
% If the project defines `latex_template`, insert content there instead.
% ============================================================

\documentclass[12pt, a4paper]{article}

% --- Encoding & fonts ---
\usepackage[T1]{fontenc}
\usepackage[utf8]{inputenc}
\usepackage{lmodern}
\usepackage{microtype}

% --- Page geometry ---
\usepackage[margin=2.5cm]{geometry}

% --- Mathematics ---
\usepackage{amsmath, amssymb, amsthm}
\usepackage{mathtools}

% Theorem environments
\newtheorem{theorem}{Theorem}[section]
\newtheorem{lemma}[theorem]{Lemma}
\newtheorem{corollary}[theorem]{Corollary}
\newtheorem{proposition}[theorem]{Proposition}
\newtheorem{definition}[theorem]{Definition}
\newtheorem{remark}[theorem]{Remark}

% --- Algorithms ---
\usepackage{algorithm}
\usepackage{algpseudocode}

% --- Code listings ---
\usepackage{listings}
\usepackage{xcolor}
\lstset{
  basicstyle=\ttfamily\small,
  breaklines=true,
  frame=single,
  numbers=left,
  numberstyle=\tiny,
  keywordstyle=\color{blue},
  commentstyle=\color{gray},
  stringstyle=\color{orange},
}

% --- Tables ---
\usepackage{booktabs}
\usepackage{tabularx}
\usepackage{multirow}

% --- Figures ---
\usepackage{graphicx}
\usepackage{subcaption}
\graphicspath{{./figures/}}

% --- Hyperlinks & cross-references ---
\usepackage[colorlinks=true, linkcolor=blue, citecolor=blue, urlcolor=blue]{hyperref}
\usepackage{cleveref}

% --- Bibliography ---
% Using biblatex with biber (preferred). Switch to natbib if the venue requires it.
\usepackage[backend=biber, style=authoryear, sorting=nyt]{biblatex}
\addbibresource{references.bib}

% --- Utilities ---
\usepackage{lipsum}   % placeholder text — remove before submission
\usepackage{todonotes}  % \todo{}, \missingfigure{} — remove before submission

% ============================================================
% Document metadata — fill in before drafting
% ============================================================

\title{%
  \textbf{Draft Title}\\
  \large Subtitle (if applicable)
}

\author{%
  First Author\thanks{Affiliation. \texttt{email@institution.edu}} \and
  Second Author\thanks{Affiliation. \texttt{email@institution.edu}}
}

\date{\today}

% ============================================================
\begin{document}
% ============================================================

\maketitle

\begin{abstract}
  One paragraph (150--250 words).
  State the problem, the approach, the key result, and the significance.
  Write the abstract last.
\end{abstract}

\tableofcontents
\newpage

% ============================================================
\section{Introduction}
\label{sec:introduction}
% ============================================================

% - Motivate the problem. Why does it matter?
% - State the gap in existing work.
% - List contributions as bullet points.
% - Describe the paper structure in one short paragraph.

\begin{itemize}
  \item \textbf{Contribution 1:} ...
  \item \textbf{Contribution 2:} ...
\end{itemize}

% ============================================================
\section{Related Work}
\label{sec:related}
% ============================================================

% Organise by theme, not chronology.
% Position this work relative to each cluster of prior work.
% \cite{} every paper discussed.

% ============================================================
\section{Background}
\label{sec:background}
% ============================================================

% Definitions, notation, and prerequisite results.
% Keep this section tight — move non-essential background to appendix.

% ============================================================
\section{Method}
\label{sec:method}
% ============================================================

% The core technical contribution.
% Include pseudocode or algorithms where the procedure is non-trivial.
% Reproduce all key equations in full with variable definitions.

\begin{equation}
  \mathcal{L}(\theta) = \ldots
  \label{eq:loss}
\end{equation}

where $\theta$ are the model parameters and \ldots

% ============================================================
\section{Experiments}
\label{sec:experiments}
% ============================================================

\subsection{Experimental Setup}

% Datasets, baselines, evaluation metrics, compute details.

\subsection{Main Results}

% Results table — reproduce key numbers; never paraphrase into prose.

\begin{table}[h]
  \centering
  \caption{Main results on [Benchmark]. Best result in \textbf{bold}, second-best \underline{underlined}.}
  \label{tab:main}
  \begin{tabular}{lcccc}
    \toprule
    \textbf{Method} & \textbf{Metric A} & \textbf{Metric B} & \textbf{Metric C} \\
    \midrule
    Baseline        & 00.0 & 00.0 & 00.0 \\
    Prior Work~\cite{ref} & 00.0 & 00.0 & 00.0 \\
    \textbf{Ours}   & \textbf{00.0} & \textbf{00.0} & \textbf{00.0} \\
    \bottomrule
  \end{tabular}
\end{table}

\subsection{Ablation Study}

% Isolate each component's contribution.

\subsection{Analysis}

% Error analysis, qualitative examples, failure modes.

% ============================================================
\section{Discussion}
\label{sec:discussion}
% ============================================================

% Interpret results. Limitations. Negative results (if any).

% ============================================================
\section{Conclusion}
\label{sec:conclusion}
% ============================================================

% Summarise contributions. Future work.

% ============================================================
% Appendix (optional)
% ============================================================

\appendix

\section{Supplementary Material}
\label{app:supplementary}

% Proofs, additional experiments, implementation details.

% ============================================================
% Bibliography
% ============================================================

\printbibliography

% ============================================================
\end{document}
% ============================================================
